\section{Further checks on maize}
\label[appendix]{app:maize-checks}

This appendix gives the details of two groups of checks summarised in \cref{sec:maize}. All were added after the main estimates and are descriptive.

\subsection{Heterogeneity of the slope}

The interaction of heat with the irrigated share of grain maize is not detectable (two-sided $p = \IrrInterP$). Nor is its interaction with the available water capacity of the soils under maize, averaged from a 90~m soil map \parencite{romandobarco2019} over the maize parcels of the RPG: it has the expected negative sign, but $p = \RuInterP$ (\RuInterBothP{} with the irrigated share in the same model, with which it correlates at $\RuCorIrr$). The 17 departments span only \RuMin--\RuMax~mm, which leaves little to work with. In the variant with the vapour pressure deficit, the interaction is detectable ($p = \RuVpdP$): the slope of the vapour pressure deficit falls from about \RuVpdLow\,\% per kPa on the shallowest soils to \RuVpdHigh\,\% on the deepest. This is one of five interaction tests on 17 values of the soil variable, so I report it as a lead to confirm, not as a finding. Department-specific slopes, shrunk towards their mean with an empirical Bayes estimator \parencite{dersimonian1986}, show little dispersion ($\tau = \EbTau$; homogeneity test $p = \EbQp$, which ignores the dependence between departments within a year and is therefore, if anything, too permissive): all shrunk slopes lie between \EbMin{} and \EbMax{}, and they correlate at $\EbCorIrr$ with the irrigated share and $+\EbCorAlt$ with the altitude gap, on 17 points. A causal forest \parencite{athey2019} is uninformative with 435 observations: its calibration test detects no heterogeneity ($p = \ForestP$) and its average effect is very imprecise.

\subsection{Composition of the land}

First, one more degree-day above 32\,$^{\circ}$C is associated with a change of $\CompEst$ in the irrigated share of grain maize area ($p = \CompP$, clustered by year, for information only). This check has little power by construction: areas are largely decided at sowing, before the summer heat. Second, and closer to the more plausible threat, rainfed maize sown for grain may be harvested as silage in dry summers. Over 2010--2025, the only years for which the SAA reports silage maize separately by irrigation mode, the share of rainfed maize harvested as silage (\SilageShare\,\% on average) does not rise detectably in hotter summers (one-sided $p = \SilageP$; added later). This check misses 2003, and silage areas may be recorded before the harvest decision is final. Abandoned fields, finally, stay in the statistics with a zero yield.
